\documentclass[11pt]{article}
\usepackage{icmlko}
\begin{document}

\title{분기는 애니의 것이지 판의 것이 아니다 --- 겹말이 아닌데도 안 는다}
\author{월드모델 실험실}
\date{노트 448 · 2026년 8월 2일}
\maketitle

\begin{abstract}
노트 447 의 계절 축은 값이 \textbf{애니 한 도메인에서} 나왔다(판
$+0.0264 \cdot 12/12$, 전이 부호까지 바꿈). 기계도 뻔했다 --- 분기 애니가
$1 \cdot 4 \cdot 7 \cdot 10$월에 몰린다. 그래서 \textbf{분기 시작까지의
거리}를 직접 적으면 같은 정보를 더 날카롭게 줄 줄 알았다.

\textbf{반증됐다.}

\begin{center}
\begin{tabular}{lrrr}
\hline
팔 & 판 & 날짜 통제 판 & 애니 \\
\hline
계절만 (챔피언) & $+0.4809$ & $+0.5199$ & $+0.5589$ \\
계절 $+$ 분기 & $+0.4812$ & $+0.5221$ & $\mathbf{+0.5756}$ \\
분기만 & $+0.4816$ & $+0.5216$ & $+0.5678$ \\
\hline
\end{tabular}
\end{center}

\textbf{애니는 $+0.0167$ 오르는데 판은 $+0.0003$ 이다.} 애니가 판의
$18\%$($606/3{,}430$)를 차지하니 산술로는 $+0.003$ 이 와야 하는데 안 온다
--- \textbf{나머지 도메인이 그만큼 돌려준다.}

노트 412 가 이미 적은 무늬다: \emph{한 도메인만 아는 열은 도메인 이름표를
한 번 더 적는 것}이다. 분기 구조는 애니에만 있고, 그것을 축으로 주면 그
축은 ``이 행은 애니다''라고 말하는 데 절반을 쓴다.
\end{abstract}

\begin{figure}[t]
\centering
\includegraphics[width=\textwidth]{figs/f1.pdf}
\caption{\textbf{가} --- 계절만 대비 변화. 초록(애니)만 서고 판은 문턱에
못 닿는다. \textbf{나} --- 계절이 이미 있던 축들과 얼마나 겹치나. 제일
붙은 것이 \texttt{trend\_momentum} 의 $-0.103$ 이고, 분기거리와도
$+0.061$ 이다 --- \textbf{겹말이 아니다.}}
\end{figure}

\section{덤 --- 계절은 달력 넷과 안 겹친다}

노트 447 이 남긴 물음 하나가 여기서 닫힌다. 이미 판에는 달력 축 넷이
있으니 계절이 그것들의 겹말일 수 있었다. 애니 학습 $2{,}779$행에서
계절과 모든 축의 스피어만을 찍어 보니 \textbf{제일 붙은 것이
$-0.103$}(\texttt{trend\_momentum})이다. \textbf{계절은 새 정보다} ---
노트 447 의 $+0.0041$ 은 있던 것을 다시 센 게 아니다.

그리고 이번 실패도 겹말 탓이 아니다. \textbf{분기거리와 계절의 상관이
$+0.061$ 밖에 안 된다} --- 둘은 거의 독립인데도 분기가 판에 아무것도
못 보탠다. \emph{정보가 겹쳐서 안 는 게 아니라, 그 정보가 한 도메인
바깥에서 쓸모가 없어서} 안 느는 것이다.

\section{판정}

\textbf{유지: ``계절 축만 있는 F18''}(노트 432 규약대로 이름으로 적는다).
등록한 문턱(판과 날짜 통제 판이 \emph{둘 다} $+0.002$ 이상)에 판이 못
닿는다. 날짜 통제 판만 $+0.0023$ 으로 문턱을 넘는데, \textbf{한쪽만
넘은 것으로는 안 넣는다} --- 노트 446 의 가드는 거부권이지 승격권이 아니다.

\section*{남은 것}

\begin{itemize}
\item \textbf{도메인 전용 축을 판에 실어 나르는 법.} 이걸로 노트 412 의
무늬가 세 번째다(시장팝업 전용 축 · 표지 둘째 · 분기). 세 번 같은 벽을
만났으면 벽을 재야 한다 --- \emph{전용 축의 값을 그 도메인 밖으로
옮기는} 장치가 있나.
\item \textbf{애니 $+0.0167$ 은 진짜다.} 판에 못 와도 애니 하나를 쓰는
쪽에는 값이 있다. 배포가 도메인별이면 쓰는 게 맞다 --- 지금 규약은
그걸 못 적는다.
\item LODO 씨앗 셋$\to$아홉 --- 아홉 노트째 미룸.
\end{itemize}

\end{document}
